\begin{table}[htbp]
\centering
\caption{Coeficientes das variáveis-chave no modelo logístico multinomial, recorte medicina, avaliação in\_sample.}
\label{tab:coeficientes_chave_logit_ternario_recorte_medicina_in_sample}
\small
\begin{tabular}{lllll}
\toprule
\rowcolor[gray]{0.85}
Especificação & Classe & Renda & Desempenho & Interação \\
\midrule
Modelo 1 & Lista de espera & 0,0949 & -3,0898 & 0,6658 \\
\rowcolor[gray]{0.95}
Modelo 1 & Não contratado & 0,0141 & 1,6764 & -0,3469 \\
Modelo 1 & Contratada & -0,1090 & 1,4134 & -0,3189 \\
\rowcolor[gray]{0.95}
Modelo 2 & Lista de espera & 0,0783 & -4,2436 & 0,3947 \\
Modelo 2 & Não contratado & 0,0090 & 2,2685 & -0,1865 \\
\rowcolor[gray]{0.95}
Modelo 2 & Contratada & -0,0873 & 1,9750 & -0,2082 \\
Modelo 3 & Lista de espera & 0,0733 & -4,8588 & 0,3267 \\
\rowcolor[gray]{0.95}
Modelo 3 & Não contratado & 0,0119 & 2,5301 & -0,1499 \\
Modelo 3 & Contratada & -0,0852 & 2,3287 & -0,1767 \\
\rowcolor[gray]{0.95}
Modelo 4 & Lista de espera & 0,0804 & -4,9498 & 0,2989 \\
Modelo 4 & Não contratado & 0,0118 & 2,5712 & -0,1423 \\
\rowcolor[gray]{0.95}
Modelo 4 & Contratada & -0,0922 & 2,3786 & -0,1566 \\
Modelo 5 & Lista de espera & 0,0395 & -5,4290 & 0,2275 \\
\rowcolor[gray]{0.95}
Modelo 5 & Não contratado & 0,0263 & 3,0019 & -0,1136 \\
Modelo 5 & Contratada & -0,0658 & 2,4271 & -0,1139 \\
\bottomrule
\end{tabular}
\par\footnotesize{Notas: Coeficientes extraídos dos modelos treinados. Em modelos multinomiais, os coeficientes pertencem às classes do desfecho; por isso, devem ser lidos como informação complementar às probabilidades previstas.}
\par\footnotesize{Fonte: elaboração própria (2026).}
\end{table}
